\chapter{Creio na ressurreição da carne}

O décimo primeiro artigo do Credo professa a fé na \textbf{ressurreição da carne}, ou seja, na ressurreição do corpo humano. Esta crença é essencial à fé cristã desde os seus inícios: ``A confiança dos cristãos é a ressurreição dos mortos; crendo nisto, vivemos''. O Catecismo da Igreja Católica, nos parágrafos 988 a 1019, explica que a ressurreição dos corpos está intimamente ligada à de Cristo, que é o primeiro ressuscitado. Todos os homens ressuscitarão com os seus corpos, mas para a vida ou para o juízo, segundo as suas obras. Esta esperança transforma a visão cristã da morte e do corpo.

\section{Introdução (988--991)}

A \textbf{ressurreição da carne} significa que, não só a alma imortal viverá após a morte, mas também o nosso corpo mortal será ressuscitado. No sepulcro semeamos um corpo corruptível, mas Deus ressuscitará um corpo incorruptível, um ``corpo espiritual'' (1 Cor 15,42-44).

A fé na ressurreição dos mortos é elemento essencial da fé cristã desde o princípio. ``Como podeis dizer que não há ressurreição dos mortos? Se não há ressurreição dos mortos, também Cristo não ressuscitou. E, se Cristo não ressuscitou, é vã a nossa pregação e vã a vossa fé'' (1 Cor 15,12-14). Na verdade, Cristo ressuscitou dos mortos, primícias dos que dormem.

\section{A ressurreição de Cristo e a nossa (992--996)}

Deus revelou progressivamente aos homens a ressurreição dos mortos. A esperança na ressurreição corporal estabeleceu-se como consequência intrínseca da fé em Deus criador do homem todo inteiro, alma e corpo. Os mártires dos Macabeus confessaram: ``O Rei do universo nos ressuscitará para uma vida eterna, porque morremos por suas leis'' (2 Mac 7,9).

Jesus ensina firmemente a ressurreição. Aos saduceus que a negavam, responde: ``Não é por isso que estais em erro, por não conhecerdes as Escrituras nem o poder de Deus?'' (Mc 12,24). A fé na ressurreição apoia-se na fé em Deus ``Deus não dos mortos, mas dos vivos'' (Mc 12,27). Jesus liga-a à sua pessoa: ``Eu sou a Ressurreição e a Vida'' (Jo 11,25).

Ser testemunha de Cristo é ser ``testemunha da sua Ressurreição''. Os encontros com o Cristo ressuscitado caracterizam a esperança cristã. Ressuscitaremos como Cristo, com Ele e por Ele.

\section{Como ressuscitarão os mortos? (997--1004)}

Na morte, a separação da alma e do corpo, o corpo decompõe-se e a alma vai ao encontro de Deus, à espera da reunião com o corpo glorificado. Deus, com o seu poder todo-poderoso, concederá definitivamente vida incorruptível aos nossos corpos, reunindo-os às nossas almas, pela virtude da Ressurreição de Jesus.

\textbf{Quem ressuscitará?} Todos os mortos: ``Os que fizeram o bem, para a ressurreição da vida, e os que fizeram o mal, para a ressurreição do julgamento'' (Jo 5,29).

\textbf{Como?} Cristo ressuscitou com o seu corpo: ``Vede as minhas mãos e os meus pés: sou eu mesmo'' (Lc 24,39); mas não voltou à vida terrena. Assim, todos ressuscitarão com os seus corpos que agora têm, mas Cristo ``transformará o nosso corpo de humilhação, conformando-o ao seu corpo glorioso'' (Fl 3,21), num ``corpo espiritual'' (1 Cor 15,44).

\textbf{Quando?} Definitivamente ``no último dia'', ``no fim do mundo''. A ressurreição dos mortos está associada à Parusia de Cristo: ``O Senhor mesmo [\ldots] descerá do céu [\ldots] e os mortos em Cristo ressuscitarão primeiro'' (1 Tes 4,16).

\section{O corpo da ressurreição e o destino da humanidade (1005--1019)}

O que é semear? A morte destrói a vida natural, mas semeia um corpo glorioso. ``O que semeias não é o corpo que há-de vir, mas um simples grão [\ldots] O que se semeia corruptível, ressuscita incorruptível [\ldots] É preciso que este corpo corruptível revista a incorruptibilidade'' (1 Cor 15,37.42.53).

Esta ``qualidade'' excede a nossa imaginação e entendimento; só é acessível à fé. A nossa participação na Eucaristia dá-nos já um antegozo da transfiguração dos nossos corpos: assim como o pão, pela bênção de Deus, deixa de ser pão comum para ser Eucaristia, assim os nossos corpos, que participam da Eucaristia, possuem a esperança da ressurreição.

Ressuscitados com Cristo, participamos já da sua vida celeste pela virtude do Espírito Santo. Pelo Batismo, fomos sepultados com Ele e ressuscitados com Ele. A vida cristã é já participação na morte e Ressurreição de Cristo.

A dignidade do corpo cristão exige respeito pelo próprio corpo e pelo dos outros, especialmente dos sofredores: ``O vosso corpo é para o Senhor [\ldots] Glorificai, pois, a Deus no vosso corpo'' (1 Cor 6,13.20).

Jesus anuncia o destino final da humanidade: uma comunhão de vida e amor com a Santíssima Trindade, no céu novo e na terra nova. Esta esperança escatológica não diminui o cuidado das realidades terrestres, mas purifica-as e confirma a ordem da criação.
